\section{Analysis}
\label{sec:analysis}

This section covers the engineering analyses that drove the major design
choices: how to handle MI-data scarcity, how to choose a base model size, how
to design an evaluation that actually informed iteration, and how to structure
safety so it stayed auditable.

\subsection{Data Scarcity and Synthetic Augmentation}

The first practical constraint was data. AnnoMI~\cite{wu2022annomi}, the only
public corpus of expert-annotated MI dialogue, contains 133 transcripts across
multiple behavioral-health domains. The alcohol-only filter reduces this to a
small slice -- after preprocessing into the project's instruction-tuning
format, only $\approx$363 training and $\approx$40 evaluation examples remain.
That is not enough material on which to train a small language model to
exhibit a consistent stylistic identity.

The choice was either to (a) train on the alcohol slice as-is and accept
overfit-style brittleness, (b) widen the source to other AnnoMI domains and
introduce off-topic drift, or (c) generate additional in-domain MI dialogue
synthetically. Option (c) was selected. The synthetic generator
(\texttt{scripts/generate\_synthetic\_data.py}) prompts Claude to simulate both
sides of a counseling session, parameterized by scenario type:

\begin{itemize}
    \item \textbf{short} -- single-turn exchange,
    \item \textbf{medium} -- 3--5 turns,
    \item \textbf{long} -- 8--12 turns,
    \item \textbf{mixed defensiveness} -- the user shifts between open and
    resistant within a single conversation.
\end{itemize}

The generator prompt embeds explicit MI guardrails: no advice-giving, no
closed questions where open ones would do, the simple-then-complex reflection
pattern, varied openers to prevent template collapse. Generation was run
in batches of $\approx$50--100 conversations and saved as JSONL under
\texttt{data/synthetic/}, then deduplicated and merged with the AnnoMI subset.

The trade-off is honest: synthetic data is generated by a model that does not
itself experience ambivalence, so the simulated client side is a stylized
imitation rather than a faithful reproduction of human ambivalence. The
evidence from evaluation (\Cref{sec:results}) is that despite this limitation,
the augmented training set substantially improves response style on real
AnnoMI test turns -- so the synthetic dialogues are at least teaching the
right output distribution.

\subsection{Model Size and Local Deployability}

A second decision was the base model. The constraint of running locally on
a CPU through Ollama eliminated 7B-and-up models for interactive latency
reasons. Among the small instruction-tuned families available at project
start, three were candidates: \texttt{Llama-3.2-1B-Instruct},
\texttt{Phi-3-mini-4k-Instruct} (3.8B), and \texttt{Gemma-2-2b-it}.
Gemma-2-2B was selected on three grounds:
\begin{enumerate}
    \item \textbf{Quantization quality.} Gemma-2 quantizes to 4-bit GGUF with
    minimal observable degradation on the project's eval set, an expected
    outcome given Gemma's training-side quantization-friendliness.
    \item \textbf{License.} Permissive enough for an academic capstone that
    redistributes a fine-tuned variant.
    \item \textbf{Tooling.} First-class support in PEFT, in TRL's SFTTrainer,
    and in Ollama as a target format.
\end{enumerate}

The 2B size yielded $\approx$1.6~GB GGUF after Q4\_K\_M quantization, with
$\sim$2--4~s per-turn latency on a CPU. Larger models would close more of the
quality gap to Claude (\Cref{sec:results}) but at a latency cost that breaks
the conversational feel.

\subsection{Evaluation Design}

The single most important methodological decision was committing to LLM-as-judge
evaluation rather than human ratings. Human MI scoring requires trained raters,
$\approx$5--10 minutes per dialogue for high-quality annotation, and inter-rater
adjudication. None of that fits the iteration loop of a one-person semester
project.

The judge is Claude Sonnet, prompted with a fixed five-dimension rubric
(empathy, MI technique, resistance handling, autonomy support, appropriateness),
each scored 1--5 with a written rationale. The rubric is reproduced in
\Cref{sec:experiments}. Three design choices follow from the recommendations
of Chiang and Lee~\cite{chiang2023closer} on LLM-judge bias:
\begin{itemize}
    \item \textbf{Rubric-based, not pairwise.} Pairwise judgements are
    susceptible to position bias and verbosity bias; absolute scoring against
    a rubric is more robust at the cost of less calibrated individual scores.
    Since the goal here is comparing configurations \emph{against the same
    rubric}, absolute scoring is sufficient.
    \item \textbf{Held-out judge.} The judge model (Claude Sonnet) is from a
    different provider than any system-under-test other than the Claude Haiku
    reference; this avoids the most pathological self-preference bias.
    \item \textbf{Coarse signal, not ground truth.} The scores are treated as
    an iteration signal. Configurations are not declared ``better'' on the
    basis of a $0.05$ difference; only effects large enough to be visible to
    a human reader are emphasized in \Cref{sec:results}.
\end{itemize}

\subsection{Safety as an Auditable Layer}

The safety design uses regex rather than an embedding classifier. The argument
for an embedding model is generalization -- it would catch crisis ideation
phrased in ways the regex misses. The argument against, in this setting, is
auditability: a regex pattern list is line-by-line inspectable, and a regression
in coverage is a one-line diff. An embedding-based classifier on the same
problem requires either a labeled crisis corpus (which the project does not
have) or an off-the-shelf classifier whose decision boundaries are not
inspectable.

The cost of the regex choice was made visible mid-project: the original
patterns required strict adjacency between intent words and self-references,
so phrasings like ``don't even want to be here'' were missed. The fix was
adding optional-intensifier groups to the patterns -- a one-commit change.
That fix would not have been a one-commit change against an opaque classifier.

The defense in depth is the audit log. Even in cases where the regex misses
a borderline phrasing, the per-turn JSONL record allows after-the-fact
discovery of the miss, which an embedding classifier would not provide.

\subsection{What This Analysis Yielded}

The analyses above produced four concrete commitments that shaped the rest of
the project:
\begin{itemize}
    \item Synthetic data is part of training, not just real AnnoMI.
    \item The base model is Gemma-2-2B, not anything larger.
    \item The eval is LLM-as-judge with a rubric, run on a mixed test set
    spanning real AnnoMI turns, synthetic scenarios, and edge cases.
    \item Safety is a two-stage regex screen with audit logging, not an
    embedding classifier.
\end{itemize}

Each of these is revisited in \Cref{sec:future_work} with the cases under
which it should be reconsidered.
